\documentclass{article}%
\usepackage[T1]{fontenc}%
\usepackage[utf8]{inputenc}%
\usepackage{lmodern}%
\usepackage{textcomp}%
\usepackage{lastpage}%
\usepackage{geometry}%
\geometry{margin=2.2cm}%
%
\usepackage{amsmath}%
\usepackage{amssymb}%
\usepackage{booktabs}%
\usepackage{longtable}%
\usepackage{array}%
\usepackage{hyperref}%
\hypersetup{hidelinks}%
\title{ITA 2026 {-} Matemática (claude\_sonnet\_5\_5) x Poliedro\\[4pt]\large Comparação com o gabarito Poliedro Resolve\thanks{Gerado com gpt-resolve: \url{https://github.com/lgabs/gpt-resolve}}}%
\date{10/11/2026}%
%
\begin{document}%
\normalsize%
\maketitle%
\section*{Resumo}\textbf{Nota média: 9,3/10} (12 de 14 itens com a resolução correta; cada item vale o mesmo dentro da questão).\par\medskip%
\begin{longtable}{@{}c c p{0.78\textwidth}@{}}\toprule Questão & Nota & Resumo \\ \midrule
1 & 10,0 & A resolução está certa e coincide com o gabarito Poliedro: 231 pontos. Ela identifica o sólido como |x|+|y|+|z| $\leq$ 5, com o argumento de vértices antípodas bem justificado, e faz a contagem por dois métodos que concordam entre si. Uma enumeração computacional confirma o total. A divisão em 129 pontos no interior e 102 na superfície não foi pedida, mas também está correta. Há só uma imprecisão de redação: a resolução fala em \textbackslash{}"losango de lado s\textbackslash{}" quando s é a semidiagonal. \\
2 & 10,0 & A resolução e o gabarito Poliedro dão a mesma resposta, n = 15, e os dois estão certos. A entrada (1,3) de A\^{}n é n(n+3)/2, conferida em Python multiplicando a matriz para n = 1..16. Ela vale exatamente 119 em n = 14, que não serve porque o enunciado pede maior que 119, e 135 em n = 15. A resolução está completa e rigorosa: usa a decomposição I + N com N nilpotente e trata corretamente a igualdade em n = 14. O gabarito tem só um erro de digitação sem importância: faltam parênteses no passo de indução. \\
3 & 10,0 & A resolução e o gabarito Poliedro dão a mesma área, 30, com os mesmos vértices: X=(3,5/2), A=(-2,0) e B=(0,4). Refiz todas as contas em Python, e elas estão corretas: as tangências, o incentro, o raio igual a Área/p, o ângulo reto em B e o produto vetorial que dá 30. A resolução segue o mesmo caminho da 2ª solução do gabarito e acrescenta a verificação de que a circunferência é a inscrita. Não encontrei erro em nenhum dos dois; só há detalhes de forma no gabarito. \\
4 & 10,0 & A resolução está correta e bate com o gabarito do Poliedro: A = $\pi$($\alpha$$^2$ + 32$\alpha$ + 528) unidades de área. Conferi o enunciado no PDF original (página 5), que coincide com a transcrição. Refiz a soma do centro com os 16 anéis vermelhos e conferi em Python que a área vermelha mais a branca dá $\pi$($\alpha$+32)$^2$, a área do disco inteiro. A resolução e o gabarito leem o enunciado da mesma forma e chegam ao mesmo resultado. Nenhum dos dois tem erro. \\
5 & 10,0 & A resolução e o gabarito do Poliedro chegam ao mesmo resultado, (22+9$\surd$6)/16 $\approx$ 2,75, com R = r(1+$\surd$6/2). Refiz a conta por três caminhos (altura do tetraedro, cubo e coordenadas) e todos confirmam. A resolução está correta e bem justificada: explica por que o centro da esfera maior é o circuncentro, confere a consistência da configuração e não comete o erro comum de tomar R igual ao circunraio. A transcrição do enunciado confere com o PDF original. Não há divergência. \\
6 & 3,3 & O gabarito e a resolução concordam que o enunciado, como foi impresso, é contraditório: as colunas n-1 e n de AB são iguais e teriam de somar 1 e 0. O Poliedro conclui que tais matrizes não existem e deixa os três itens sem resposta. A resolução conclui a mesma coisa, mas responde como principal uma leitura corrigida, trocando as colunas de AB pelas de A: det(AB) = 0, tr(AB) = n(n-1)/2 e, para n = 5, $\alpha$1 = 1, $\alpha$2 $\in$ \{1,-2\}, $\alpha$3 = 1, $\alpha$4 livre. Conferi essas contas e elas estão certas nessa leitura. No item a), det(AB) = 0 vale para qualquer escolha dos $\alpha$, então nenhum dos dois erra. Nos itens b) e c), pelo enunciado como foi impresso, o gabarito está certo. Os valores da resolução respondem a uma versão alterada da questão, embora a própria resolução registre que o conjunto solução literal é vazio. A nota final depende de como a banca tratou a questão (anulação ou leitura aceita), o que deve ser verificado. \\
7 & 10,0 & A resolução e o gabarito do Poliedro chegam ao mesmo p(x) = 2x\^{}4 {-} 3x\^{}3 + 7x\^{}2 + 7x {-} 5 e às mesmas raízes: {-}1, 1/2, 1+2i e 1{-}2i. Conferi as contas de forma independente e estão certas. A resolução está completa: justifica por que o quociente tem grau 1, verifica o resto e confere a soma e o produto das raízes no final. Não há pontos duvidosos. \\
8 & 10,0 & A resolução e o gabarito Poliedro dão a mesma resposta, 784, e os dois estão certos. A resolução troca a, b, c pelo primeiro termo e pelas diferenças (a, d, e) e soma separando a par e a ímpar (408 + 376). O Poliedro parte de C(99,2), tira as soluções com duas coordenadas iguais e divide por 3! = 6. Conferi por enumeração direta em Python, que também dá 784, e pela fórmula das partições em três partes distintas. Nenhum dos dois tem erro, e a solução está completa e bem justificada. \\
9 & 10,0 & A resolução está correta e bate com o Poliedro nos dois itens. a) A hipérbole é x$^2$/4 - y$^2$/32 = 1 e as tangentes por C são y = $\pm$2$\surd$6$\cdot$x - 8. b) Os pontos de tangência são ($\surd$6, 4) para m = 2$\surd$6 e (-$\surd$6, 4) para m = -2$\surd$6. A resolução ainda justifica por que descarta a reta vertical e as paralelas às assíntotas, e confere os resultados pela condição de tangência e pela polar y = 4. A única falha está no gabarito: na frase final ele troca os rótulos T$_{1}$/T$_{2}$ em relação às retas t$_{1}$/t$_{2}$, mas no desenvolvimento a associação está certa. \\
10 & 10,0 & A resolução e o gabarito do Poliedro chegam às mesmas respostas nos dois itens, e as duas estão corretas: (a) r\_C = 2$\surd$3 cm e r\_D = 6 cm; (b) S = 10$\pi$ - 12$\surd$3 cm$^2$ $\approx$ 10,63 cm$^2$. Conferi o enunciado no PDF original e refiz tudo em Python: AP = PQ = 6, BP = AP$\cdot$$\surd$(2+$\surd$3), CP = CA, DP = DB, e a área bate com a fórmula geral da interseção de círculos. A resolução também justifica que M fica entre A e B e que os dois segmentos circulares são menores. Só há problemas de redação no parágrafo de verificação do item b (uma expressão truncada e a estimativa arredondada para 10,2 em vez de 10,14), sem efeito nas respostas. \\
\bottomrule\end{longtable}%
\section*{Principais pontos de dúvida}%
\begin{itemize}%
\item%
\textbf{Q6} \textit{(statement)} Lido ao pé da letra, o enunciado é impossível para todo n $\geq$ 2. A coluna n de A é nula, então as colunas n-1 e n de AB são iguais, e a hipótese exige que a mesma soma valha 1 e 0. Para n $\geq$ 3 há contradição até se ignorarmos a coluna n, porque s\_\{n-1\} = a\_\{n,n-1\} = 1, enquanto as diferenças das somas acumuladas exigem -1. Vale verificar se o ITA anulou a questão ou publicou errata. É isso que decide se a resposta esperada é 'não existem matrizes' ou alguma leitura corrigida.%
\item%
\textbf{Q6} \textit{(solution)} As respostas principais da resolução em b) (n(n-1)/2) e em c) ($\alpha$1 = 1, $\alpha$2 $\in$ \{1,-2\}, $\alpha$3 = 1, $\alpha$4 livre) dependem de trocar 'colunas de AB' por 'colunas de A'. A própria resolução admite que essa leitura é uma conjectura. As contas estão certas dentro dessa leitura, e ela é plausível: com $\alpha$\_j = 1, cada coluna de A soma exatamente n-j, e o traço sai em função de n. Mesmo assim, se a correção seguir o enunciado impresso, essas respostas não valem. A conclusão literal, 'contraditório, conjunto vazio', aparece só como ressalva.%
\end{itemize}%
\section*{Questão 1}%
\href{https://poliedroresolve.sistemapoliedro.com.br/vestibulares/ita/2026/ita-2-fase-dia-1-28-10-2025-2-fase-dia-1/dissertativa-1-matematica-2-fase-dia-1-ita-2026}{Resolução do gabarito}\par\smallskip%
\begin{longtable}{@{}p{0.06\textwidth} p{0.33\textwidth} p{0.33\textwidth} p{0.1\textwidth} p{0.08\textwidth}@{}}\toprule Item & Resolução & Gabarito & Veredito & Certo \\ \midrule
Questão 1 -- número de pontos de coordenadas inteiras no octaedro (interior + superfície) & 231 pontos (129 no interior e 102 na superfície) & 231 pontos & igual & ambos \\
\bottomrule\end{longtable}%
\begin{itemize}%
\item%
\textbf{Questão 1 -- número de pontos de coordenadas inteiras no octaedro (interior + superfície))} Releio o enunciado no PDF original e ele confere com a transcrição. Os vértices são ($\pm$5,0,0), (0,$\pm$5,0) e (0,0,$\pm$5), e o sólido é |x|+|y|+|z| $\leq$ 5. A resolução conta por s=|x|+|y|: 11 + $\Sigma$ 4s(11-2s) = 11 + 220 = 231. Também confere por camadas: 1+6+18+38+66+102 = 231. O gabarito fatia por z: 2(1+5+13+25+41)+61 = 231. Fiz a enumeração em Python e ela deu 231, sendo 102 pontos com |x|+|y|+|z| = 5 (superfície) e 129 no interior. A divisão interior/superfície da resolução não foi pedida, mas está correta. A fórmula (2n+1)(2n$^2$+2n+3)/3 dá 7, 25 e 231 para n = 1, 2 e 5, o que também confere.%
\end{itemize}%
\paragraph{Pontos de dúvida}%
\begin{itemize}%
\item%
\textit{(solution)} A resolução chama o conjunto |x|+|y|=s de "losango de lado s". Na verdade o lado mede s$\surd$2, e s é a semidiagonal. A contagem de 4s pontos está certa, só o termo está impreciso.%
\item%
\textit{(statement)} O texto diz "no interior e na superfície". A leitura natural, adotada pela resolução e pelo gabarito, é a união: o total de pontos do sólido fechado. Uma leitura literal como interseção seria vazia e não faz sentido, então não há ambiguidade real.%
\end{itemize}%
\section*{Questão 2}%
\href{https://poliedroresolve.sistemapoliedro.com.br/vestibulares/ita/2026/ita-2-fase-dia-1-28-10-2025-2-fase-dia-1/dissertativa-2-matematica-2-fase-dia-1-ita-2026}{Resolução do gabarito}\par\smallskip%
\begin{longtable}{@{}p{0.06\textwidth} p{0.33\textwidth} p{0.33\textwidth} p{0.1\textwidth} p{0.08\textwidth}@{}}\toprule Item & Resolução & Gabarito & Veredito & Certo \\ \midrule
Q2 -- menor n $\in$ $\mathbb{N}$ com (A\^{}n)\_\{13\} > 119 & n = 15 ((A\^{}n)\_\{13\} = n(n+3)/2, que vale 119 em n=14 e 135 em n=15) & n = 15 ((A\^{}n)\_\{13\} = (n$^2$+3n)/2 > 119 $\Rightarrow$ n > 14) & igual & ambos \\
\bottomrule\end{longtable}%
\begin{itemize}%
\item%
\textbf{Q2 -- menor n $\in$ $\mathbb{N}$ com (A\^{}n)\_\{13\} > 119)} O enunciado do PDF (p. 3) confere com a transcrição. Fórmula conferida em Python, multiplicando a matriz diretamente para n = 1..16: (A\^{}n)\_\{13\} = n(n+3)/2 em todos os casos, com 104 (n=13), 119 (n=14) e 135 (n=15). Como a desigualdade é estrita, n=14 dá exatamente 119 e não serve, então n=15. A resolução usa A = I + N com N nilpotente (N$^3$ = 0) e o binômio de Newton; o gabarito chega à mesma fórmula por indução. Os dois métodos são válidos. A inequação (n-14)(n+17) > 0 dá n > 14 ou n < -17, como diz o gabarito.%
\end{itemize}%
\paragraph{Pontos de dúvida}%
\begin{itemize}%
\item%
\textit{(key)} Erro de digitação no passo de indução do gabarito: a entrada (1,3) de A\^{}\{k+1\} aparece como '(k + 1)\^{}2 + 3(k + 1)/2', sem os parênteses. O correto é {[}(k+1)$^2$ + 3(k+1){]}/2. Também não aparece a conta intermediária (k$^2$+3k)/2 + 2k + 2 = {[}(k+1)$^2$ + 3(k+1){]}/2. Isso não muda a resposta.%
\end{itemize}%
\section*{Questão 3}%
\href{https://poliedroresolve.sistemapoliedro.com.br/vestibulares/ita/2026/ita-2-fase-dia-1-28-10-2025-2-fase-dia-1/dissertativa-3-matematica-2-fase-dia-1-ita-2026}{Resolução do gabarito}\par\smallskip%
\begin{longtable}{@{}p{0.06\textwidth} p{0.33\textwidth} p{0.33\textwidth} p{0.1\textwidth} p{0.08\textwidth}@{}}\toprule Item & Resolução & Gabarito & Veredito & Certo \\ \midrule
Área de ABCD & 30 (X=(3,5/2), A=(-2,0), B=(0,4), D=(6,1); {[}ABX{]}=15/2; {[}ABCD{]}=4$\cdot${[}ABX{]}) & 30 (X=(3,5/2), A=(-2,0), B=(0,4); ABX retângulo em B, {[}ABX{]}=15/2; {[}ABCD{]}=4$\cdot${[}ABX{]}) & igual & ambos \\
\bottomrule\end{longtable}%
\begin{itemize}%
\item%
\textbf{Área de ABCD)} Conferi o enunciado no PDF original (p. 4): a transcrição está fiel. Refiz as contas em Python. X=(3,5/2). C pertence a r, então A=2X-C=(-2,0). As tangentes por A à circunferência de centro (1/2,5/2) e raio $\surd$5/2 têm m=1/2 (que é a própria r) e m=2, logo AB: y=2x+4 e B=AB?s=(0,4), D=(6,1). O incentro ponderado dá (1/2, 5/2) e r=Área/p=(15/2)/(3$\surd$5)=$\surd$5/2. Portanto a circunferência é a inscrita, e não uma exinscrita. (A-B)$\cdot$(X-B)=0, o que confirma o ângulo reto em B. {[}ABX{]}=15/2 e {[}ABCD{]}=|AB$\times$AD|=30. A resolução segue o mesmo caminho da 2ª solução do gabarito e ainda verifica que a circunferência é a inscrita. Nas conferências deste trabalho, a resolução está um pouco mais completa que o gabarito.%
\end{itemize}%
\paragraph{Pontos de dúvida}%
\begin{itemize}%
\item%
\textit{(key)} Na 1ª solução, o gabarito escreve AB = 4$\surd$5/2 sem simplificar (o valor é 2$\surd$5). O valor está correto e é só questão de forma. Ele também afirma 'Da figura, IX é paralela ao eixo x', o que de fato vale, porque I e X têm y=5/2, mas não justifica pelas coordenadas.%
\item%
\textit{(statement)} O enunciado não diz qual vértice fica em s, B ou D. O que define B=(0,4), e não D=(6,1), é a circunferência inscrita em ABX. As duas soluções usam isso. A resolução mostra por que a outra tangente (m=1/2, que é a própria r) é descartada e confirma pelo incentro que se trata da circunferência inscrita, e não de uma exinscrita.%
\end{itemize}%
\section*{Questão 4}%
\href{https://poliedroresolve.sistemapoliedro.com.br/vestibulares/ita/2026/ita-2-fase-dia-1-28-10-2025-2-fase-dia-1/dissertativa-4-matematica-2-fase-dia-1-ita-2026}{Resolução do gabarito}\par\smallskip%
\begin{longtable}{@{}p{0.06\textwidth} p{0.33\textwidth} p{0.33\textwidth} p{0.1\textwidth} p{0.08\textwidth}@{}}\toprule Item & Resolução & Gabarito & Veredito & Certo \\ \midrule
Q4 -- área total em vermelho em função de $\alpha$ & A\_verm = $\pi$($\alpha$$^2$ + 32$\alpha$ + 528) unidades de área & S = $\pi$($\alpha$$^2$ + 32$\alpha$ + 528) unidades de área & igual & ambos \\
\bottomrule\end{longtable}%
\begin{itemize}%
\item%
\textbf{Q4 -- área total em vermelho em função de $\alpha$)} Refiz a conta. Os 33 círculos têm raios $\alpha$, $\alpha$+1, \ldots{}, $\alpha$+32. O círculo central é vermelho; o anel entre $\alpha$+k-1 e $\alpha$+k tem área $\pi$(2$\alpha$+2k-1) e é vermelho quando k é par (k = 2, 4, \ldots{}, 32). São 16 anéis vermelhos, cuja soma dá $\pi${[}16(2$\alpha$-1) + 2$\cdot$272{]} = $\pi$(32$\alpha$ + 528). Somando o centro, a área vermelha é $\pi$($\alpha$$^2$ + 32$\alpha$ + 528). Conferi em Python. Os anéis brancos somam $\pi$(32$\alpha$ + 496), e vermelho mais branco dá $\pi$($\alpha$+32)$^2$, a área do disco inteiro. A resolução organiza a soma por k = 2j, e o gabarito soma a PA 2$\alpha$+3, 2$\alpha$+7, \ldots{}, 2$\alpha$+63. É a mesma soma, com o mesmo resultado.%
\end{itemize}%
\paragraph{Pontos de dúvida}%
\begin{itemize}%
\item%
\textit{(statement)} A resolução e o gabarito leem "33 círculos" da mesma forma: o círculo central conta como o primeiro, e os raios vão de $\alpha$ a $\alpha$+32, formando o centro e 32 anéis. Também supõem que o primeiro anel, logo depois do centro vermelho, é branco, o que decorre de "cores alternadas". Por isso o último anel é vermelho. Essa é a leitura natural do enunciado. Uma leitura com 33 círculos além do central mudaria a resposta, mas é pouco plausível.%
\item%
\textit{(solution)} A verificação com $\alpha$ = 0 ("sem círculo central") não corresponde ao enunciado, que pressupõe um círculo central de raio $\alpha$ > 0. Serve apenas como checagem algébrica da fórmula e não afeta a resposta.%
\end{itemize}%
\section*{Questão 5}%
\href{https://poliedroresolve.sistemapoliedro.com.br/vestibulares/ita/2026/ita-2-fase-dia-1-28-10-2025-2-fase-dia-1/dissertativa-5-matematica-2-fase-dia-1-ita-2026}{Resolução do gabarito}\par\smallskip%
\begin{longtable}{@{}p{0.06\textwidth} p{0.33\textwidth} p{0.33\textwidth} p{0.1\textwidth} p{0.08\textwidth}@{}}\toprule Item & Resolução & Gabarito & Veredito & Certo \\ \midrule
Questão 5 (razão V\_maior / soma dos volumes das 4 menores) & (22+9$\surd$6)/16 $\approx$ 2,753, com R = r(2+$\surd$6)/2 & (1/4)(1+$\surd$6/2)$^3$ = (22+9$\surd$6)/16, com R = r(1+$\surd$6/2) & igual & ambos \\
\bottomrule\end{longtable}%
\begin{itemize}%
\item%
\textbf{Questão 5 (razão V\_maior / soma dos volumes das 4 menores))} Os dois dão a mesma resposta final, (22+9$\surd$6)/16, e o mesmo raio: R = r(2+$\surd$6)/2 = r(1+$\surd$6/2). Refiz a conta. A aresta do tetraedro é 2r e a altura é h = 2r$\surd$6/3. O circunraio vale 3h/4 = r$\surd$6/2, logo R = r + r$\surd$6/2. Depois, (2+$\surd$6)$^3$ = 44+18$\surd$6, e (44+18$\surd$6)/32 = (22+9$\surd$6)/16 $\approx$ 2,7528. A conferência por coordenadas da resolução também está certa: centros em ($\pm$1,$\pm$1,$\pm$1) com número par de sinais negativos, aresta 2$\surd$2, r = $\surd$2 e R = $\surd$3+$\surd$2. Isso dá ($\surd$3+$\surd$2)$^3$ = 9$\surd$3+11$\surd$2, e dividindo por 8$\surd$2 sai de novo (22+9$\surd$6)/16. Conferi tudo em Python. Os métodos diferem só na forma: a resolução acha o circunraio pelo cubo e por Pitágoras, e o gabarito usa a razão 3:1 na altura. O enunciado transcrito confere com o PDF original (página 6).%
\end{itemize}%
\paragraph{Pontos de dúvida}%
\begin{itemize}%
\item%
\textit{(key)} Na transcrição do gabarito, a expressão "R = r + 3/42r(sqrt(6))/3" ficou embaralhada. A leitura correta é R = r + (3/4)$\cdot$(2r$\surd$6/3) = r(1+$\surd$6/2). É um problema de formatação, sem efeito no resultado.%
\item%
\textit{(statement)} O enunciado diz "inscritas" sem detalhar. A resolução declara a hipótese de tangência interna de cada esfera menor à maior, e o gabarito a usa implicitamente. As duas leituras coincidem, e é a interpretação padrão.%
\end{itemize}%
\section*{Questão 6}%
\href{https://poliedroresolve.sistemapoliedro.com.br/vestibulares/ita/2026/ita-2-fase-dia-1-28-10-2025-2-fase-dia-1/dissertativa-6-matematica-2-fase-dia-1-ita-2026}{Resolução do gabarito}\par\smallskip%
\begin{longtable}{@{}p{0.06\textwidth} p{0.33\textwidth} p{0.33\textwidth} p{0.1\textwidth} p{0.08\textwidth}@{}}\toprule Item & Resolução & Gabarito & Veredito & Certo \\ \midrule
a & det(AB) = 0, para qualquer n $\geq$ 2 e quaisquer $\alpha$\_j, porque a linha 1 de A é nula. A resolução também diz que a hipótese literal sobre as colunas de AB é impossível. & Sem resposta: não existem matrizes com as propriedades descritas. As diferenças entre as somas acumuladas dão S\_i = -1 para i > 1, o que contradiz S\_n = 0. & diferente & ambos \\
b & tr(AB) = n(n-1)/2, na leitura 'soma da coluna j de A = n-j'. Pela hipótese literal sobre AB, a resolução observa que o traço teria de valer 1 e 0 ao mesmo tempo e que nenhum valor fica determinado. & Sem resposta: o enunciado é inconsistente e não existem tais matrizes. & diferente & gabarito \\
c & $\alpha$1 = 1, $\alpha$2 $\in$ \{1, -2\}, $\alpha$3 = 1 e $\alpha$4 $\in$ $\mathbb{R}$ qualquer, na leitura sobre as colunas de A. A resolução acrescenta que, com o enunciado literal, o conjunto solução seria vazio. & Não há valores possíveis para $\alpha$1, $\alpha$2, $\alpha$3, $\alpha$4, porque não existem tais matrizes. & diferente & gabarito \\
\bottomrule\end{longtable}%
\begin{itemize}%
\item%
\textbf{a)} Conferi com frações exatas para n = 2 a 7 e $\alpha$ aleatórios: as colunas n-1 e n de AB são sempre iguais, porque a coluna n de A é nula. Com isso, a condição 'soma da coluna j de AB = n-j' exige que a mesma soma valha 1 e 0. O enunciado literal é contraditório e o Poliedro está certo ao dizer que tais matrizes não existem. Por outro lado, det(AB) = 0 vale para todas as matrizes com essa definição, já que a linha 1 de A e de AB é nula e det B = 1. Então a afirmação da resolução é verdadeira de forma incondicional (sob a hipótese, vale por vacuidade). Ela também aponta a mesma inconsistência que o gabarito. Nenhum dos dois erra a matemática: o gabarito recusa responder e a resolução dá o valor que vale para qualquer escolha dos $\alpha$\_j.%
\item%
\textbf{b)} Pelo enunciado como foi impresso, o original prevalece. A hipótese nunca é satisfeita para n $\geq$ 2, então não há traço a determinar, e a conclusão do Poliedro está certa. A resolução chega à mesma conclusão sobre o enunciado literal, mas coloca como resposta principal o valor obtido trocando a hipótese de AB para A. Essa conta está certa. Conferi tr(AB) = $\Sigma$ a\_il b\_li = soma de todos os elementos de A, porque a\_il $\neq$ 0 só quando l < i e aí b\_li = 1. Com s\_j = n-j isso dá $\Sigma$(n-j) = n(n-1)/2, e o caso n = 5 dá 10. Só que essa é a resposta de um enunciado modificado por conjectura da própria resolução, e não do enunciado da prova.%
\item%
\textbf{c)} Pelo enunciado literal, nenhuma quádrupla serve e o conjunto solução é vazio, como diz o Poliedro. A própria resolução registra isso na resposta final, mas como ressalva, e destaca os valores da leitura alternativa. Conferi esses valores: 1+x+x$^2$+x$^3$ = 4 $\Leftrightarrow$ (x-1)(x$^2$+2x+3) = 0, com $\Delta$ = -8, logo x = 1. Também 1+x+x$^2$ = 3 $\Leftrightarrow$ x $\in$ \{1, -2\} e 1+x = 2 $\Leftrightarrow$ x = 1. Como $\alpha$4 só aparece como $\alpha$4$^{0}$, ele fica livre. Com (1,-2,1,3) e (1,1,1,7), as colunas de A somam 4,3,2,1,0. As colunas de AB, porém, somam 4,7,9,10,10, o que confirma que esses $\alpha$ não satisfazem a hipótese impressa.%
\end{itemize}%
\paragraph{Pontos de dúvida}%
\begin{itemize}%
\item%
\textit{(statement)} Lido ao pé da letra, o enunciado é impossível para todo n $\geq$ 2. A coluna n de A é nula, então as colunas n-1 e n de AB são iguais, e a hipótese exige que a mesma soma valha 1 e 0. Para n $\geq$ 3 há contradição até se ignorarmos a coluna n, porque s\_\{n-1\} = a\_\{n,n-1\} = 1, enquanto as diferenças das somas acumuladas exigem -1. Vale verificar se o ITA anulou a questão ou publicou errata. É isso que decide se a resposta esperada é 'não existem matrizes' ou alguma leitura corrigida.%
\item%
\textit{(solution)} As respostas principais da resolução em b) (n(n-1)/2) e em c) ($\alpha$1 = 1, $\alpha$2 $\in$ \{1,-2\}, $\alpha$3 = 1, $\alpha$4 livre) dependem de trocar 'colunas de AB' por 'colunas de A'. A própria resolução admite que essa leitura é uma conjectura. As contas estão certas dentro dessa leitura, e ela é plausível: com $\alpha$\_j = 1, cada coluna de A soma exatamente n-j, e o traço sai em função de n. Mesmo assim, se a correção seguir o enunciado impresso, essas respostas não valem. A conclusão literal, 'contraditório, conjunto vazio', aparece só como ressalva.%
\item%
\textit{(key)} A resolução do Poliedro é sucinta, mas está correta: S\_1 + \ldots{} + S\_k = n-k dá S\_k = -1 para k $\geq$ 2, o que contradiz S\_n = 0 (e também S\_\{n-1\} = 1). Há uma imprecisão de notação na matriz C: a última coluna aparece como $\Sigma$\_\{i=1\}\^{}j $\alpha$\_i\^{}\{n-i-1\} com um j solto. Na verdade a coluna n de C é igual à coluna n-1. Isso não muda a conclusão.%
\item%
\textit{(key)} No item a), o determinante é 0 para qualquer escolha dos $\alpha$\_j, porque a linha 1 de AB é nula. Mesmo com a hipótese inconsistente, um corretor pode aceitar 'det(AB) = 0'. Deixar o item sem resposta, como faz o gabarito, é rigoroso, mas não é a única saída razoável.%
\item%
\textit{(statement)} Há duas hipóteses implícitas. Primeira: é preciso adotar a convenção 0$^{0}$ = 1 para a\_\{j+1,j\} = $\alpha$\_j$^{0}$ quando $\alpha$\_j = 0. A resolução declara essa convenção; o gabarito não. Segunda: o caso trivial n = 1 (AB = {[}0{]}) satisfaz a hipótese, então a inexistência vale só para n $\geq$ 2. Nenhum dos dois comenta esse caso, e ele não afeta c), em que n = 5.%
\end{itemize}%
\section*{Questão 7}%
\href{https://poliedroresolve.sistemapoliedro.com.br/vestibulares/ita/2026/ita-2-fase-dia-1-28-10-2025-2-fase-dia-1/dissertativa-7-matematica-2-fase-dia-1-ita-2026}{Resolução do gabarito}\par\smallskip%
\begin{longtable}{@{}p{0.06\textwidth} p{0.33\textwidth} p{0.33\textwidth} p{0.1\textwidth} p{0.08\textwidth}@{}}\toprule Item & Resolução & Gabarito & Veredito & Certo \\ \midrule
Raízes de p(x) em $\mathbb{C}$ (e o polinômio p(x)) & p(x) = 2x\^{}4 {-} 3x\^{}3 + 7x\^{}2 + 7x {-} 5 = (x+1)(2x{-}1)(x\^{}2{-}2x+5); raízes: {-}1, 1/2, 1+2i, 1{-}2i & p(x) = 2x\^{}4 {-} 3x\^{}3 + 7x\^{}2 + 7x {-} 5; raízes: {-}1, 1/2, 1 + 2i, 1 {-} 2i & igual & ambos \\
\bottomrule\end{longtable}%
\begin{itemize}%
\item%
\textbf{Raízes de p(x) em $\mathbb{C}$ (e o polinômio p(x)))} Conferi tudo de novo em Python, com frações exatas. A divisão de 2x\^{}4{-}3x\^{}3+7x\^{}2+7x{-}5 por x\^{}3{-}x+4 dá quociente 2x{-}3 e resto 9x\^{}2{-}4x+7. A soma dos coeficientes que entram em Girard dá {-}b/a = 3/2 e o produto dá e/a = {-}5/2. Substituindo, p({-}1) = p(1/2) = p(1$\pm$2i) = 0. A resolução e o gabarito seguem o mesmo caminho: escrever p = ($\alpha$x+$\beta$)(x\^{}3{-}x+4) + resto, aplicar Girard ($\alpha$=2, $\beta$={-}3) e depois Briot{-}Ruffini por {-}1 e por 1/2. O que sobra é 2x\^{}2{-}4x+10 (ou x\^{}2{-}2x+5), cujas raízes são 1$\pm$2i. O enunciado do PDF original (p. 8) é igual à transcrição.%
\end{itemize}%
\section*{Questão 8}%
\href{https://poliedroresolve.sistemapoliedro.com.br/vestibulares/ita/2026/ita-2-fase-dia-1-28-10-2025-2-fase-dia-1/dissertativa-8-matematica-2-fase-dia-1-ita-2026}{Resolução do gabarito}\par\smallskip%
\begin{longtable}{@{}p{0.06\textwidth} p{0.33\textwidth} p{0.33\textwidth} p{0.1\textwidth} p{0.08\textwidth}@{}}\toprule Item & Resolução & Gabarito & Veredito & Certo \\ \midrule
Questão 8 (número de sequências estritamente crescentes de três naturais com soma 100) & 784 & 784 & igual & ambos \\
\bottomrule\end{longtable}%
\begin{itemize}%
\item%
\textbf{Questão 8 (número de sequências estritamente crescentes de três naturais com soma 100))} Conferi o enunciado no PDF original (p. 9) e ele bate com a transcrição. A convenção da prova é $\mathbb{N}$ = \{1, 2, 3, \ldots{}\}, e tanto a resolução quanto o gabarito usam a, b, c $\geq$ 1.\newline%
\newline%
Resolução: faz a = a, d = b - a $\geq$ 1, e = c - b $\geq$ 1 e chega a 3a + 2d + e = 100. Isso equivale a contar os pares (a, d) com 3a + 2d $\leq$ 99. Para a ímpar há $\Sigma$(51 - 3k) = 408 pares, e para a par há $\Sigma$(49 - 3k) = 376, com k = 1..16. Os limites estão certos: a vai até 32, e para k = 16 restam 3 e 1 valores de d. Total: 784.\newline%
\newline%
Gabarito: faz a = x + 1, b = y + 1, c = z + 1. A equação x + y + z = 97 tem C(99,2) = 4851 soluções não negativas. Os três casos com duas coordenadas iguais têm 49 soluções cada (x = 0..48). Não existe x = y = z, porque 3 não divide 97. Dividindo por 6: (4851 - 147)/6 = 784.\newline%
\newline%
Conferi por computador: a enumeração direta dá 784, e as duas fórmulas e as somas parciais 408/376 também. A verificação da resolução pelo inteiro mais próximo de (n - 3)$^2$/12 também vale: 97$^2$/12 $\approx$ 784,08.%
\end{itemize}%
\paragraph{Pontos de dúvida}%
\begin{itemize}%
\item%
\textit{(statement)} A resposta depende de $\mathbb{N}$ excluir o zero. Essa convenção está escrita explicitamente na prova ($\mathbb{N}$ = \{1, 2, 3, \ldots{}\}), então não há ambiguidade. Se o zero fosse permitido, entrariam também as 49 ternas (0, b, 100 - b) com 1 $\leq$ b $\leq$ 49, e o total seria 833. A resolução e o gabarito usam a convenção correta.%
\end{itemize}%
\section*{Questão 9}%
\href{https://poliedroresolve.sistemapoliedro.com.br/vestibulares/ita/2026/ita-2-fase-dia-1-28-10-2025-2-fase-dia-1/dissertativa-9-matematica-2-fase-dia-1-ita-2026}{Resolução do gabarito}\par\smallskip%
\begin{longtable}{@{}p{0.06\textwidth} p{0.33\textwidth} p{0.33\textwidth} p{0.1\textwidth} p{0.08\textwidth}@{}}\toprule Item & Resolução & Gabarito & Veredito & Certo \\ \midrule
a & H: x$^2$/4 - y$^2$/32 = 1; tangentes por C=(0,-8): y = 2$\surd$6$\cdot$x - 8 e y = -2$\surd$6$\cdot$x - 8 & H: x$^2$/4 - y$^2$/32 = 1 (8x$^2$ - y$^2$ = 32); t$_{1}$: y = 2$\surd$6$\cdot$x - 8 e t$_{2}$: y = -2$\surd$6$\cdot$x - 8 & igual & ambos \\
b & y = 2$\surd$6$\cdot$x - 8 tangente em ($\surd$6, 4); y = -2$\surd$6$\cdot$x - 8 tangente em (-$\surd$6, 4) & Para t$_{1}$ (m=2$\surd$6): ($\surd$6, 4); para t$_{2}$ (m=-2$\surd$6): (-$\surd$6, 4). Na frase final o gabarito escreve T$_{1}$=(-$\surd$6,4) e T$_{2}$=($\surd$6,4), com os rótulos trocados & igual & ambos \\
\bottomrule\end{longtable}%
\begin{itemize}%
\item%
\textbf{a)} Refiz a conta: c=6, e=3 $\Rightarrow$ a=2, b$^2$=32. Substituindo y=mx-8 chega{-}se a (8-m$^2$)x$^2$+16mx-96=0, com $\Delta$=3072-128m$^2$=0 $\Rightarrow$ m=$\pm$2$\surd$6. A resolução ainda descarta com razão a reta x=0, que não encontra H porque exige |x|$\geq$2, e as paralelas às assíntotas (m$^2$=8), que cortam H em um só ponto sem ser tangentes. Também confere o resultado pela condição n$^2$=a$^2$m$^2$-b$^2$.%
\item%
\textbf{b)} Com $\Delta$=0, a raiz dupla é x\_T = -16m/(2(8-m$^2$)) = m/2. Para m=2$\surd$6, x=$\surd$6 e y=4; para m=-2$\surd$6, x=-$\surd$6 e y=4. Os dois pontos satisfazem 8$\cdot$6-16=32, e a conferência numérica deu o mesmo. Estão também sobre a polar de C, y=4. A associação reta?ponto da resolução bate com o desenvolvimento do gabarito. A única diferença é a frase final do gabarito, que troca os rótulos T$_{1}$/T$_{2}$ em relação a t$_{1}$/t$_{2}$. Isso é um deslize de digitação do gabarito e não afeta a resolução.%
\end{itemize}%
\paragraph{Pontos de dúvida}%
\begin{itemize}%
\item%
\textit{(key)} Na conclusão do item b o gabarito escreve T$_{1}$=(-$\surd$6,4) e T$_{2}$=($\surd$6,4), mas no desenvolvimento ele associa m=2$\surd$6 (t$_{1}$) a ($\surd$6,4) e m=-2$\surd$6 (t$_{2}$) a (-$\surd$6,4). Quem ler T\_i como o ponto da reta t\_i encontra o par trocado. O conjunto de pontos e a associação no corpo da resolução estão certos, e a resolução conferida usa a associação correta.%
\item%
\textit{(key)} O gabarito decide o sinal de y (y>0) olhando a figura, sem justificar algebricamente. Isso não afeta o resultado: o próprio y=mx-8 com x=m/2 dá y=4. A resolução conferida obtém os pontos sem depender da figura.%
\end{itemize}%
\section*{Questão 10}%
\href{https://poliedroresolve.sistemapoliedro.com.br/vestibulares/ita/2026/ita-2-fase-dia-1-28-10-2025-2-fase-dia-1/dissertativa-10-matematica-2-fase-dia-1-ita-2026}{Resolução do gabarito}\par\smallskip%
\begin{longtable}{@{}p{0.06\textwidth} p{0.33\textwidth} p{0.33\textwidth} p{0.1\textwidth} p{0.08\textwidth}@{}}\toprule Item & Resolução & Gabarito & Veredito & Certo \\ \midrule
a & r\_C = 2$\surd$3 cm; r\_D = 6 cm (com AP = PQ = 6 cm) & CP = 2$\surd$3 cm; DP = 6 cm & igual & ambos \\
b & S = (10$\pi$ - 12$\surd$3) cm$^2$ $\approx$ 10,63 cm$^2$ & S = 10$\pi$ - 12$\surd$3 cm$^2$ & igual & ambos \\
\bottomrule\end{longtable}%
\begin{itemize}%
\item%
\textbf{a)} Refiz a conta. Com A na origem: AP = PQ = 6, P = (3$\surd$3, 3), C em 2$\surd$3, D em 6$\surd$3, B em 6+6$\surd$3. Daí CP = 2$\surd$3 = CA, DP = 6 = DB e BP = 6$\surd$(2+$\surd$3) $\approx$ 11,591, o que bate com o enunciado. Os caminhos são diferentes. A resolução acha x = AP pela soma AM + MB = AB e depois tira os raios. O gabarito usa o ângulo PBM = 15$^\circ$ e o ângulo central PDM = 30$^\circ$, chega a k = $\surd$3 e daí aos mesmos raios. Os dois estão corretos.%
\item%
\textbf{b)} A lente é a soma de dois segmentos circulares menores. O de raio 2$\surd$3 e ângulo central 120$^\circ$ tem área 4$\pi$ - 3$\surd$3, e o de raio 6 e ângulo 60$^\circ$ tem área 6$\pi$ - 9$\surd$3. Somando, 10$\pi$ - 12$\surd$3 $\approx$ 10,6313 cm$^2$. Conferi pela fórmula geral da interseção de dois círculos (d = CD = 4$\surd$3) e deu 10,631317, o mesmo valor.%
\end{itemize}%
\paragraph{Pontos de dúvida}%
\begin{itemize}%
\item%
\textit{(solution)} No parágrafo 'Verificação' do item b, a expressão '(C+r\_1){-}(D{-}r\_2)=4\textbackslash{}sqrt3{-}(3\textbackslash{}sqrt3{-}6)\textbackslash{}cdot0+\textbackslash{}dots' está truncada e não faz sentido. A largura certa da lente é 4$\surd$3 - (6$\surd$3 - 6) = 6 - 2$\surd$3 $\approx$ 2,536 cm. A estimativa (2/3)$\cdot$2,536$\cdot$6 dá $\approx$ 10,14, mas o texto diz $\approx$ 10,2. Isso não afeta a resposta, mas fica mal redigido numa prova discursiva.%
\item%
\textit{(solution)} No cálculo de r\_2, a resolução começa por t = (MB$^2$ - 9)/(2MB) e muda de rumo no meio ('mas é mais direto usar...'). A conta final r\_2 = BP$^2$/(2MB) = 6 está correta, mas o trecho ficou confuso.%
\end{itemize}%
\end{document}